\section{Measure complete mixing calls}\label{pcm-measurement}
Each measurement includes function entry, clipping-count initialization, complete
groups, scalar tail and return. It excludes caller preparation, assembly, session
creation, WAV writing and program exit. Both implementations return the same samples
and clipping counts. The scalar version uses explicit bounds; the MSA version adds
eight-sample groups with a population-count reduction.

\begin{table}[htbp]
\centering
\caption{Native instructions for quiet and positive-clipping inputs. Samples count mono component values in one function call.}
\tagpdfsetup{table/header-rows={1}}
\begin{tabular}{rrrrr}
\toprule
Samples & Scalar quiet & MSA quiet & Scalar clipped & MSA clipped\\
\midrule
0 & 6 & 9 & 6 & 9\\
1 & 36 & 39 & 34 & 37\\
7 & 216 & 219 & 202 & 205\\
8 & 246 & 29 & 230 & 29\\
9 & 276 & 59 & 258 & 57\\
16 & 486 & 49 & 454 & 49\\
63 & 1896 & 359 & 1770 & 345\\
64 & 1926 & 169 & 1798 & 169\\
\bottomrule
\end{tabular}
\end{table}

The quiet inputs have small opposite-signed values with no clipping. The positive
clipping inputs pair 32767 with one at every position. Scalar branch paths differ
between these families. Full vector groups execute the same instruction sequence
for either family, while a scalar tail retains the data-dependent branch cost.
Short calls pay for the vector eligibility test without executing a full group.

A 63-sample call uses seven groups and seven scalar samples. A 64-sample call uses
eight complete groups. This explains why the extra sample can reduce native work.
At 64 samples, a mixed positive/negative clipping pattern uses 1894 scalar or 169
MSA instructions. Replacing either input exactly in place gives the same counts
for these implementations because both input groups are loaded before each store.

Both paths observe four source bytes read and two destination bytes written per
sample. Vector loads and stores publish per-byte memory notices in the pinned core;
they do not identify host bus transactions. Unlike the previous chapter's overlapping
windows, these groups introduce no repeated transfer bytes. The work counts describe
these native implementations and inputs, not host elapsed time, and they are not
exact grading requirements for an equivalent solution.
